\color{unidarkgreen}{\subsection{Токовая защита нулевой последовательности (ТЗНП)}}
\color{black}

\begin{enumerate}

%%%%%%%%%%%%%%%%%%%%%%%%%%%%%%%%%%%%%%%%%%%%%%%%%%%%%%%%%%%%%%%%%%%%%%  ОБЩИЕ СВЕДЕНИЯ %%%%%%%%%%%%%%%%%%%%%%%%%%%%%%%%%%%%%%%%%%%%%%%%%%
\item Общие сведения
\begin{enumerate}

\item
Токовая защита нулевой последовательности (ТЗНП) предназначена для селективного отключения поврежденного участка электрической сети при коротких замыканиях на землю, реагируя на повышение тока нулевой последовательности.

\item
В устройстве реализовано шесть ступеней ТЗНП, ввод в работу которых осуществляется программными ключами \textbf{«ТЗНП-N»}, где N – номер ступени ТЗНП.

\item
Ток срабатывания ступеней ТЗНП задается с помощью уставок \textbf{«3I0ср ТЗНП-N»}. Действие всех ступеней осуществляется с регулируемыми независимыми выдержками времени \textbf{«Тср ТЗНП-N»}.

\item
На рисунке \ref{pict_tznp:TZNP} приведена функциональная схема ТЗНП-1, выполнение функциональных схем ступеней ТЗНП-2(3,4,5,6) аналогичное.

\begin{figure}[H]
\centering
\includegraphics[width=0.9\textwidth,height=0.9\textheight,keepaspectratio]{\fbpath/03.01 ЛВ/561.0111 ТЗНП/img/TZNP.pdf}
\captionof{figure}{Функциональная схема ступени ТЗНП-1}\label{pict_tznp:TZNP}
\end{figure}

\item
Уставки ступеней ТЗНП представлены в таблице \ref{app_set:tznp}.

\item
Имеется возможность задать различные режимы работы ступеней ТЗНП при возникновении неисправности в цепях напряжения. Вариант действия логики ступени ТЗНП при срабатывании БНН задается программным переключателем \textbf{«БНН ТЗНП-N»}, имеющим три положения:

\begin{itemize}
	\item \textit{Выведено} -- режим работы ступени не зависит от БНН;
	\item \textit{Вывод напр.} -- ступень переводится в ненаправленный режим работы;
	\item \textit{Вывод ст.} --  при срабатывании БНН ступень выводится из работы (однако, если ступень задана ненаправленной, то она останется в работе независимо от срабатывания БНН; ступень также не будет выводиться из работы, если направленность автоматически вывелась при срабатывании ТЗНП или в процессе АУ).
\end{itemize}

\item
Предусмотрен автоматический вывод направленности автоматически ускоряемой ступени ТЗНП (вводится программным ключом \textbf{«Вывод напр. при АУ»}). При наличии разрешающего сигнала ввода ускорения при включении \textit{«АУ: Ввод ускорения»} направленность ТЗНП выводится. Причем сигнал вывода направленности подхватывается от сигнала срабатывания АУ ТЗНП, что обеспечивает устойчивое состояние срабатывания ТЗНП при неполнофазном включении выключателя. 

\item
Кроме того, направленность может выводиться при срабатывании ТЗНП. Это обеспечивает устойчивое срабатывание ТЗНП при неполнофазном отключении, что необходимо для надежного действия УРОВ. Данный режим вводится программным ключом \textbf{«Вывод напр. при ср. ТЗНП»}.

\item
Для обеспечения несрабатывания ТЗНП от бросков намагничивающего тока предусмотрена блокировка ступеней ТЗНП по содержанию второй гармоники в токе нулевой последовательности. 
Блокировка основана на контроле отношения действующего значения второй гармоники тока нулевой последовательности к основной гармонике. Ступень ТЗНП блокируется при превышении данного соотношения пороговой величины, заданной уставкой \textbf{«Кблок 2 гарм.»}. Наличие функции блокировки от бросков намагничивающего тока задается для каждой ступени отдельно с помощью программного ключа \textbf{«БНТ ТЗНП-N»}. Для предотвращения излишней блокировки при больших токах КЗ, сопровождающихся насыщением ТТ, блокировка выводится из действия при превышении тока $3I_0$ основной гармоники уставки \textbf{«3I0блок ОВ БНТ»}.

\item
Функциональная схема логики органа выявления БНТ приведена на рисунке \ref{pict_tznp:bnt}. Уставки органа выявления БНТ представлены в таблице \ref{app_set:tznp}.

\begin{figure}[H]
\centering
\includegraphics[width=0.65\textwidth,height=0.65\textheight,keepaspectratio]{\fbpath/03.01 ЛВ/561.0111 ТЗНП/img/bnt.pdf}
\captionof{figure}{Функциональная схема логики органа выявления БНТ}\label{pict_tznp:bnt}
\end{figure}

\item
Для предотвращения неправильной работы защиты предусмотрена блокировка при протекании через трансформаторы тока двух выключателей сквозного тока внешнего КЗ, вызывающего неодновременное насыщение трансформаторов. Блокировка вводится отдельными программными ключами для первой ступени ТЗНП (\textbf{«БСТО ТЗНП-1»}), оперативного ускорения (\textbf{«БСТО ОУ ТЗНП»}).

\item
ТЗНП может быть оперативно выведена из работы от внешнего сигнала \textit{«ТЗНП: Вывод»}, назначенного на соответствующий дискретный вход или функциональную клавишу. Выводятся все ступени ТЗНП, а также все виды ускорения ТЗНП.

\item
Предусмотрена блокировка чувствительных ступеней ТЗНП, который производится от внешнего сигнала \textit{«Логика БЧС: Блок. ЧС ТЗНП»}, поданного на соответствующий вход ФСУ. Программными ключами \textbf{«Блок ТЗНП-N по чувств.»} можно задать ступени, выводимые из работы.


\end{enumerate}
%%%%%%%%%%%%%%%%%%%%%%%%%%%%%%%%%%%%%%%%%%%%%%%%%%%%%%%%%%%%%%%%%%%%%%  КОНЕЦ ОБЩИЕ СВЕДЕНИЯ %%%%%%%%%%%%%%%%%%%%%%%%%%%%%%%%%%%%%%%%%%%%%%%%%%

%%%%%%%%%%%%%%%%%%%%%%%%%%%%%%%%%%%%%%%%%%%%%%%%%%%%%%%%%%%%%%%%%%%%%% ОНМ %%%%%%%%%%%%%%%%%%%%%%%%%%%%%%%%%%%%%%%%%%%%%%%%%%
\item Органы направления мощности
\begin{enumerate}

\item
Направленность ТЗНП обеспечивается органами направления мощности нулевой последовательности (ОНМНП) и обратной последовательности (ОНМОП). 

\item
ОНМНП реагирует на величины векторов тока и напряжения нулевой последовательности, а также угол сдвига между ними. 

\item
Так как мощность нулевой последовательности всегда направлена от места несимметрии во внешнюю сеть, то разрешающий орган направления мощности (ОНМНП-Р) срабатывает при направлении мощности нулевой последовательности от линии к шинам (при КЗ впереди), а блокирующий (ОНМНП-Б) – при обратном направлении мощности. 

\item
За положительное направление отсчета векторов токов и напряжений принято направление против часовой стрелки. Углы максимальной чувствительности откладываются от вектора напряжения $\vec {3U_0}$ (вектор поляризации) и задаются уставками \textbf{«Фмч ОНМНП-Р»} и \textbf{«Фмч ОНМНП-Б»}. Уставки по углам максимальной чувствительности рекомендуется задавать для ОНМНП-Р – плюс $110^\circ$, а для ОНМНП-Б – минус $70^\circ$. Ширина зоны срабатывания $\pm 85^\circ$ от угла максимальной чувствительности.

\item
Пороги срабатывания ОНМНП-Р и ОНМНП-Б по току и напряжению регулируются уставками \textbf{«3I0ср ОНМНП-Р»}, \textbf{«3U0ср ОНМНП-Р»}, \textbf{«3I0ср ОНМНП-Б»}, и \textbf{«3U0ср ОНМНП-Б»}. Выбор напряжения нулевой последовательности, по которому работает ОНМНП, осуществляется с помощью программного переключателя \textbf{«Выбор 3U0»}.

\item
Для повышения чувствительности ОНМНП прямой направленности (ОНМНП-Р) предусмотрена возможность «искусственного смещения точки подключения ТН в линию». Напряжение нулевой 
последовательности, подводимое к ОНМНП-Р, рассчитывается с учетом добавочного напряжения смещения по формуле:

\begin{equation}
\label{eq:3U}
3\vec U_{0\textsl{ ОНМНП}} = 3\vec U_{0\textsl{ КЗ}} + 3\vec U_{0\textsl{ см}} = 3\vec U_{0\textsl{ КЗ}} + 3 \vec I_{0 \textsl{ КЗ}} \cdot \vec{Z}_{0\textsl{ см}}
\end{equation}
\begin{eqexpl}[10mm]
\item{$3\vec U_{0\textsl{ КЗ}}$}  измеряемое устройством напряжение нулевой последовательности в момент КЗ;
\item{$3\vec I_{0 \textsl{ КЗ}}$}  измеряемый устройством ток нулевой последовательности в момент КЗ;
\item{$3\vec U_{0\textsl{ см}}$}  напряжение смещения нулевой последовательности;
\item{$\vec{Z}_{0\textsl{ см}}$}  сопротивление смещения в линию (по нулевой последовательности).
\end{eqexpl}

\item
Для смещения точки подключения ТН с шин в линию аргумент сопротивления смещения необходимо принять на $180^\circ$ больше угла максимальной чувствительности. Сопротивление смещения задается уставками \textbf{«Rсм ОНМНП-Р»} и \textbf{«Хсм ОНМНП-Р»}.

\item
ОНМОП реагирует на величины векторов тока и напряжения обратной последовательности, а также угол сдвига между ними.

\item
Разрешающий орган (ОНМОП-Р) срабатывает при направлении мощности обратной последовательности от линии к шинам, а блокирующий (ОНМОП-Б) -- при обратном направлении мощности. 

\item
Углы максимальной чувствительности задаются уставками \textbf{«Фмч ОНМОП-Р»} и \textbf{«Фмч ОНМОП-Б»}. Уставки по углам максимальной чувствительности рекомендуется задавать для ОНМОП-Р -- плюс $110^\circ$, а для ОНМОП-Б -- минус $70^\circ$. Ширина зоны срабатывания $\pm 85^\circ$ от угла максимальной чувствительности.

\item
Пороги срабатывания ОНМОП-Р и ОНМОП-Б по току и напряжению регулируются уставками \textbf{«3I0ср ОНМОП-Р»}, \textbf{«3U0ср ОНМОП-Р»}, \textbf{«3I0ср ОНМОП-Б»}, и \textbf{«3U0ср ОНМОП-Б»}.

\item
Параметры ОНМНП(ОНМОП) приведены в таблице \ref{desc_tznp:tbl2}.

\small
\setlength{\extrarowheight}{0.05cm} % высота ячеек в таблице
\begin{longtable}
{|>{\centering\arraybackslash}m{15.2cm}|>{\centering\arraybackslash}m{1.7cm}|}
\caption{Параметры ОНМНП(ОНМОП)\hfill\vspace{-0.5\baselineskip}}\label{desc_tznp:tbl2}\\ 
\hhline{|-|-|}
\rowcolor{gray!30}
Наименование параметра & Значение  \\ 
\hhline{|-|-|}
\endfirsthead

\caption*{\hspace{3pt}\emph{Продолжение таблицы \ref{desc_tznp:tbl2}\hfill\vspace{-0.5\baselineskip}}} \\ % сделано по ГОСТ 2.105 п.6.8.7
\hhline{|-|-|}
\rowcolor{gray!30}
Наименование параметра & Значение  \\ 
\hhline{|-|-|}
\endhead
\hhline{|-|-|}

\endfoot
\endlastfoot

\raggedright Время срабатывания ОНМ при утроенных по отношению к порогам срабатывания значениях тока и напряжения, мс, не более & \centering\arraybackslash 35  \\  \hline
\raggedright Время возврата ОНМ при одновременном сбросе входных величин тока и напряжения от номинальных значений до нуля, мс, не более  & \centering\arraybackslash 25  \\  \hline
\raggedright Средняя основная погрешность порогов срабатывания ОНМ по току и напряжению, \% уставки, не более   & \centering\arraybackslash $\pm$5  \\  \hline
\raggedright Средняя основная абсолютная погрешность ОНМ по углу Фмч, эл. град., не более & \centering\arraybackslash $\pm$5 \\  \hline
\raggedright Дополнительная погрешность параметров срабатывания от изменения температуры окружающего воздуха в рабочем диапазоне, \% от среднего значения параметров, измеренных при температуре $(20\pm5)^\circ$C, не более & \centering\arraybackslash $\pm$5 \\  \hline
\raggedright Минимальная угловая ширина зоны срабатывания ОНМ, эл. град, не менее & \centering\arraybackslash 165 \\ \hline
\raggedright Коэффициент возврата ОНМ НП по напряжению и току, не менее  & \centering\arraybackslash 0,9\\  \hline

\end{longtable}
\setlength{\extrarowheight}{0.0cm} % возврат параметра высота ячеек в таблице
\normalsize

\item
На рисунке \ref{pict_tznp:char_triggering_ONMNP} представлены характеристики срабатывания ОНМНП-Р и ОНМНП-Б. Аналогичные характеристики для ОНМОП-Р и ОНМОП-Б.

\begin{figure}[H]
\centering
\includegraphics[width=1\textwidth,height=1\textheight,keepaspectratio]{\fbpath/03.01 ЛВ/561.0111 ТЗНП/img/char_triggering_ONMNP.pdf}
\captionof{figure}{Характеристики срабатывания ОНМНП}\label{pict_tznp:char_triggering_ONMNP}
\end{figure}

\item
Для выбора режима направленности по мощности нулевой последовательности для каждой ступени ТЗНП предусмотрен программный переключатель \textbf{«Направл. ТЗНП-N»}, имеющий три положения:

\begin{itemize}
    \item \textit{Ненапр.} -- ступень выполнена ненаправленной;
    \item \textit{Прямо} -- ступень выполнена прямонаправленной;
    \item \textit{Обратно} -- ступень выполнена обратнонаправленной.
\end{itemize}

\item
Кроме того, с помощью программного переключателя \textbf{«Контр. направл. ТЗНП-N»} можно установить различные алгоритмы контроля направленности.

\begin{itemize}
    \item \textit{Разр.} --  прямая направленность обеспечивается срабатыванием ОНМНП-Р, обратная -- срабатыванием ОНМНП-Б;
    \item \textit{Блок.} -- прямая направленность обеспечивается отсутствием срабатывания ОНМНП-Б, обратная -- отсутствием срабатывания ОНМНП-Р;
    \item \textit{Разр. или блок.} -- прямая направленность обеспечивается срабатыванием ОНМНП-Р или отсутствием срабатывания ОНМНП-Б, обратная -- срабатыванием ОНМНП-Б или отсутствием срабатывания ОНМНП-Р.
\end{itemize}

\item
Положение переключателя \textit{«Разр. или блок.»} рекомендуется выставлять для ступеней ТЗНП, выполняющих функцию дальнего резервирования, так как при удаленных КЗ напряжение $3U_0$ может оказаться недостаточным для срабатывания ОНМ из-за его уменьшения по мере удаления от точки КЗ. Уставки для выбора режима направленности по мощности обратной последовательности аналогичные.

\item
Тип ОНМ, используемый для определения направленности ступени ТЗНП, выбирается программным переключателем \textbf{«Тип ОНМ ТЗНП-N»}, имеющим четыре положения:

\begin{itemize}
    \item \textit{НП} -- для определения направленности ступени учитывается только ОНМНП;
    \item \textit{ОП} -- для определения направленности ступени учитывается только ОНМОП;
    \item \textit{НП или ОП} -- для определения направленности ступени учитывается один из органов, ОНМНП или ОНМОП;
    \item \textit{НП и ОП} -- для определения направленности ступени одновременно учитываются оба органа, ОНМНП и ОНМОП.
\end{itemize}

\item
Функциональная схема логики ОНМНП приведена на рисунке \ref{pict_tznp:onnp}.

\item
Функциональная схема логики ОНМОП приведена на рисунке \ref{pict_tznp:onop}.

\item
Уставки ОНМОП и ОНМНП представлены в таблице \ref{app_set:tznp}.

\medskip

\begin{figure}[H]
\centering
\includegraphics[width=0.8\textwidth,height=0.8\textheight,keepaspectratio]{\fbpath/03.01 ЛВ/561.0111 ТЗНП/img/onnp.pdf}
\captionof{figure}{Функциональная схема логики ОНМНП}\label{pict_tznp:onnp}
\end{figure}

\medskip

\begin{figure}[H]
\centering
\includegraphics[width=0.8\textwidth,height=0.8\textheight,keepaspectratio]{\fbpath/03.01 ЛВ/561.0111 ТЗНП/img/onop.pdf}
\captionof{figure}{Функциональная схема логики ОНМОП}\label{pict_tznp:onop}
\end{figure}

\medskip

\end{enumerate}

%%%%%%%%%%%%%%%%%%%%%%%%%%%%%%%%%%%%%%%%%%%%%%%%%%%%%%%%%%%%%%%%%%%%%% АУ %%%%%%%%%%%%%%%%%%%%%%%%%%%%%%%%%%%%%%%%%%%%%%%%%%
\item Автоматическое ускорение (АУ)
\begin{enumerate}

\item
В составе ТЗНП предусмотрена логика формирования пуска автоматического ускорения ТЗНП (АУ ТЗНП), выходной сигнал которого действует в логику автоматического ускорения. Формирование сигнала пуска АУ ТЗНП можно задать от любой ступени ТЗНП программным переключателем \textbf{«Ступень АУ ТЗНП»}. Сигнал пуска АУ ТЗНП формируется непосредственно от сигналов пуска ступени ТЗНП, то есть настройка автоматически ускоряемой ступени определяется соответствующими настройками самой ступени.

\item
Функциональная схема логики АУ ТЗНП приведена на рисунке \ref{pict_tznp:AU_tznp}.

\vspace{2mm}
\begin{figure}[H]
\centering
\includegraphics[width=1\textwidth,height=1\textheight,keepaspectratio]{\fbpath/03.01 ЛВ/561.0111 ТЗНП/img/AU_tznp.pdf}
\captionof{figure}{Функциональная схема пуска АУ ТЗНП}\label{pict_tznp:AU_tznp}
\end{figure}

\item
Уставки АУ ТЗНП представлены в таблице \ref{app_set:tznp}.


\end{enumerate}

%%%%%%%%%%%%%%%%%%%%%%%%%%%%%%%%%%%%%%%%%%%%%%%%%%%%%%%%%%%%%%%%%%%%%% ОУ %%%%%%%%%%%%%%%%%%%%%%%%%%%%%%%%%%%%%%%%%%%%%%%%%%
\item Оперативное ускорение (ОУ)
\begin{enumerate}

\item
В устройстве реализовано два логических узла оперативного ускорения ТЗНП: ОУ1 ТЗНП и ОУ2 ТЗНП. Ввод в работу ОУ ТЗНП осуществляется от внешнего сигнала \textit{«ОУ ТЗНП: Ввод»}. Оперативное ускорение ТЗНП осуществляется от ИО ТЗНП той ступени, которая выбрана программным переключателем \textbf{«Ступень ОУ ТЗНП»}. Время срабатывания ОУ ТЗНП задается уставкой \textbf{«Тср ОУ ТЗНП»}.

\item
Направленность, режим работы при БНН настраиваются как у обычной ступени.

\item
Предусмотрена возможность автоматического ввода ОУ ТЗНП при приеме внешнего сигнала о выводе основной защиты. Автоматический ввод ОУ ТЗНП разрешается программным ключом \textbf{«Ввод ОУ ТЗНП при выводе осн. защиты»}. Введенное ускорение может быть оперативно выведено от внешнего сигнала \textit{«ОУ ТЗНП: Вывод авт. введенного ОУ»}.

\item
Функциональная схема логики ОУ1 ТЗНП приведена на рисунке \ref{pict_tznp:OU_tznp}. Функциональная схема логики ОУ2 ТЗНП выполнена аналогично ОУ1 ТЗНП.

\begin{figure}[H]
\centering
\includegraphics[width=1\textwidth,height=1\textheight,keepaspectratio]{\fbpath/03.01 ЛВ/561.0111 ТЗНП/img/OU_tznp.pdf}
\captionof{figure}{Функциональная схема ОУ ТЗНП}\label{pict_tznp:OU_tznp}
\end{figure}

\medskip

\item
Уставки ОУ ТЗНП представлены в таблице \ref{app_set:tznp}.

\end{enumerate}

%%%%%%%%%%%%%%%%%%%%%%%%%%%%%%%%%%%%%%%%%%%%%%%%%%%%%%%%%%%%%%%%%%%%%% ПУ %%%%%%%%%%%%%%%%%%%%%%%%%%%%%%%%%%%%%%%%%%%%%%%%%%
\item Поперечное ускорение (ПУ)
\begin{enumerate}

\item
В устройстве предусмотрено поперечное ускорение ТЗНП (ПУ ТЗНП) от защит параллельной линии, принцип действия которого основан на сравнении направления мощности нулевой последовательности защищаемой и параллельной линии.

\item
Программный переключатель \textbf{«Ступень ПУ ТЗНП»} задает ускоряемую ступень ТЗНП. Поперечное ускорение осуществляется непосредственно от сигналов пуска ступени ТЗНП, то есть настройка ускоряемой ступени определяется соответствующими настройками самой ступени.

\item
Функциональная схема логики ПУ ТЗНП приведена на рисунке \ref{pict_tznp:PU_tznp}.

\vspace{2mm}
\begin{figure}[H]
\centering
\includegraphics[width=1\textwidth,height=1\textheight,keepaspectratio]{\fbpath/03.01 ЛВ/561.0111 ТЗНП/img/PU_tznp.pdf}
\captionof{figure}{Функциональная схема логики пуска ПУ ТЗНП}\label{pict_tznp:PU_tznp}
\end{figure}

\item
Уставки ПУ ТЗНП представлены в таблице \ref{app_set:tznp}.

\item
Данный вид ускорения может корректно работать только при наличии связи между параллельными линиями. Если шины соединены через ШСВ, то необходим контроль его включенного состояния (сигнал \textit{«ПУ ТЗНП: ШСВ/СВ включен»}). Если ШСВ выведен в ремонт или отсутствует, но имеется связь между параллельными линиями с помощью оперативного переключателя подается дискретный сигнал \textit{«ПУ ТЗНП: ШСВ/СВ выведен»}.

\item
Если ШСВ выведен, то для формирования сигнала срабатывания ПУ ТЗНП достаточно приёма сигнала \textit{«ПУ ТЗНП: ПУ от пар. ЛЭП»} (сигнал срабатывания ОНМНП-Б и РПВ параллельной линии), что свидетельствует об отсутствии повреждения на параллельной линии. Использование сигнала РПВ параллельной ВЛ позволяет исключить неправильное действие защиты при повреждении на параллельной ВЛ в зоне между выносными трансформаторами тока и одним из выключателей этой линии.

\item
Если ШСВ в работе, то необходимо контролировать его включенное положение, в данном случае для формирования сигнала срабатывания ПУ ТЗНП необходимо наличие двух сигналов: \textit{«ПУ ТЗНП: ПУ от пар. ЛЭП»} и \textit{«ПУ ТЗНП: ШСВ/CВ включен»}.

\item
Время срабатывания поперечного ускорения ТЗНП определяется уставкой \textbf{«Тср ПУ ТЗНП»}.

\item
Поперечное ускорение ТЗНП может быть выведено из работы от внешнего сигнала \textit{«ПУ ТЗНП: Вывод»}, назначенного на соответствующий дискретный вход или функциональную клавишу.

\item
Формирование сигнала \textit{«ПУ ТЗНП: Разр. ПУ пар. ЛЭП»} предназначено для выдачи на устройство РЗА, установленное на параллельной ЛЭП. Он то и должен назначаться на вход ФСУ \textit{«ПУ ТЗНП: ПУ от пар. ЛЭП»} устройства на параллельной линии. 

\end{enumerate}

\end{enumerate}